% =============================================================
%  ch03_odometrie.tex  --  Odometrie und Sensorfusion
% =============================================================

\chapter{Odometrie und Sensorfusion}

\textbf{Odometrie} schätzt die Eigenbewegung aus inkrementellen Sensordaten.
\begin{itemize}
\item \textbf{Radodometrie:} Aus Encoder-Ticks und
  Differentialantriebs-Kinematik die Pose integrieren. Einfach, aber driftet
  (Schlupf, Rundungsfehler akkumulieren).
\item \textbf{Visuelle Odometrie (VO):} Bewegung aus der Verfolgung von
  Bildmerkmalen über aufeinanderfolgende Frames schätzen. Mit Stereokamera
  bekommst du den Maßstab gratis (Tiefe ist bekannt).
\item \textbf{Visual-Inertial Odometry (VIO):} VO + IMU per EKF fusioniert --
  robuster, glatter.
\end{itemize}
In ROS\,2 ist \textbf{\code{robot\_localization}} das Standardpaket: ein
\code{ekf\_node}/\code{ukf\_node} fusioniert mehrere Quellen (Rad-Odom + IMU +
visuelle Odom) zu einer konsistenten \code{odom}$\rightarrow$\code{base\_link}-Transformation.
Drift wird später durch SLAM/Loop-Closure korrigiert
(\code{map}$\rightarrow$\code{odom}).
